\section*{Chapitre \ref{ch:rs:universe-symbology-and-corporate-actions} --- Univers, symbologie et opérations sur titres}

\begin{solution}{exo:rs:universe-symbology-and-corporate-actions:1}
$12/3 = 4$~: le cours ajusté est de 4~USD, donc un plancher sur les cours ajustés exclut une action qui s'échangeait à 12~USD, parce qu'elle
a été divisée par la suite (après une hausse).
\end{solution}

\begin{solution}{exo:rs:universe-symbology-and-corporate-actions:2}
$24/0.80 - 1 = 29$, soit un rendement de $+2\,900~\%$ entre deux sociétés différentes.
\end{solution}

\begin{solution}{exo:rs:universe-symbology-and-corporate-actions:3}
Chaque mois, 11.5~\% du capital est vendu et 11.5~\% acheté, à 0.2~\% chacun~: $12 \times 0.115 \times 2 \times 0.002 = 0.55~\%$ du capital par
an. Avec le tampon, $12 \times 0.031 \times 0.004 = 0.15~\%$.
\end{solution}

\begin{solution}{exo:rs:universe-symbology-and-corporate-actions:4}
$0.02 \times 0.30 = 0.6$ points par an~: le \omterm{def:rs:universe-symbology-and-corporate-actions:delisting}{rendement de radiation} manquant est perdu précisément sur les positions qui échouent.
\end{solution}

\begin{solution}{exo:rs:universe-symbology-and-corporate-actions:5}
Les rangs proches de la frontière bougent de quelques places par mois dans les deux sens~; une valeur au rang 499 a autant de chances d'être au rang 502 le mois suivant, comme
au rang 496, et sans tampon, chaque franchissement de ce type est un remplacement. Avec un rang de sortie de 600, un membre doit perdre une centaine de places pour
sortir, ce qui suppose un vrai changement de sa liquidité. Les membres conservés restent dans les 600 premiers, si bien que la liquidité de l'univers est presque
inchangée, tandis que les remplacements inutiles disparaissent.
\end{solution}

\begin{solution}{exo:rs:universe-symbology-and-corporate-actions:6}
La série ajustée des divisions~: la distance est un rapport de deux cours dans la fenêtre, et les facteurs postérieurs s'y annulent, tandis que les divisions situées
dans la fenêtre sont retirées. Les cours échangés transforment une division située dans la fenêtre en une baisse apparente de 50~\% ou plus~; une série de rendement total
intègre les dividendes dans un niveau de cours, si bien qu'une action à fort dividende paraît plus éloignée de son plus haut que ne l'est son cours.
\end{solution}

\begin{solution}{exo:rs:universe-symbology-and-corporate-actions:7}
Avec une probabilité de 0, ni ticker réutilisé ni rendement raccordé~; avec 0.35, 164 rendements raccordés (158 tickers réutilisés), de taille absolue moyenne
164~\% et de médiane 71~\%~; avec 0.7, 398 (371 tickers), de moyenne 259~\% et de médiane 90~\%. Le nombre croît avec la probabilité~; chaque taille est le
rapport de deux cours sans lien, et la moyenne est dominée par les rares cas où le dernier cours de l'ancien détenteur était proche de zéro (une
radiation pour performance), que la médiane ignore.
\end{solution}

\begin{solution}{exo:rs:universe-symbology-and-corporate-actions:8}
Le fichier actuel ne contient que les survivants d'aujourd'hui (\omterm{def:rs:point-in-time-data-and-the-biases:survivorship}{biais de survie}), et ses cours ajustés reportent dans le passé les divisions ultérieures (le
plancher écarte les futurs gagnants)~; reconstruire l'historique jusqu'en 2010 à partir de lui applique les deux biais à chaque date. L'univers doit provenir d'un
\omterm{def:rs:universe-symbology-and-corporate-actions:master}{référentiel des titres} (security master) à la date (point-in-time) incluant les valeurs radiées, du plancher sur les cours échangés, et de l'exigence d'historique comptée depuis chaque date
à rebours.
\end{solution}

\begin{solution}{pb:rs:universe-symbology-and-corporate-actions:1}
\textbf{1.} 1\,471 cotations, dont 471 prennent fin. \textbf{2.} 113 pour performance, 358 par fusion. \textbf{3.} 335 divisions dans 231 valeurs.
\textbf{4.} 158 tickers sur 1\,549. \textbf{5.} 164. \textbf{6.} 164~\% (médiane 71~\%). \textbf{7.} $39.55/4.23 - 1 = +835~\%$.
\textbf{8.} 0.30~\%. \textbf{9.} Les fenêtres raccordées sont seulement celles qui chevauchent un changement, et l'univers liquide détient rarement les
sociétés défaillantes avant leur radiation ni les nouvelles cotations peu après la leur. \textbf{10.} 0.34~\%. \textbf{11.} 71~\% contre
15~\%. \textbf{12.} Une action est divisée parce que son cours a monté~; le cours ajusté divise les cours anciens par la division ultérieure, si bien que les
valeurs repoussées sous le plancher sont celles qui ont ensuite monté. \textbf{13.} 8.65 échangé et 4.32 ajusté au mois~0~; une division de deux pour un
division au mois~57. \textbf{14.} Toute règle portant sur un niveau~: un seuil de capitalisation calculé avec les cours ajustés et les actions actuelles, un
cours minimal pour les ventes à découvert, une distance à un plus haut passé calculée avec les cours de rendement total. \textbf{15.} 11.5~\% et 3.1~\% par mois.
\textbf{16.} Presque rien~: chaque membre figure dans les 600 premiers par valeur échangée. \textbf{17.} La simulation inscrit le
rendement du motif au mois de radiation ($-30\%$ ou le prix de la fusion)~; dans les données réelles, on prend le registre des radiations du \omterm{def:rs:universe-symbology-and-corporate-actions:master}{référentiel des titres}, avec une
estimation par motif lorsque le rendement manque. \textbf{18.} \emph{Résultat nommé~:} une fois jointe sur les tickers, la table de cours simulée
contient 164 rendements couvrant deux sociétés, de taille absolue moyenne de 164~\%~; un plancher de 5~USD sur les cours ajustés retire 0.34~\% des
mois-titres dont le rendement de l'année suivante est de 71~\%, contre 15~\% pour l'univers. \textbf{19.} Toute table indexée par ticker, et toute
liste de composants «~actuels~» utilisée pour le passé. \textbf{20.} Un ticker ne désigne une cotation que sur une plage de dates, et la clé ne doit pas
changer de sens au fil du temps.
\end{solution}

\begin{solution}{iq:rs:universe-symbology-and-corporate-actions:1}
Les tickers changent (renommages) et sont réutilisés après les radiations, si bien qu'un ticker peut désigner plusieurs sociétés au fil du temps et une société
plusieurs tickers. Indexer tout sur un identifiant interne permanent, avec une \omterm{def:rs:universe-symbology-and-corporate-actions:permanent}{correspondance des identifiants} de chaque symbole externe vers cet identifiant, datée par plage de dates.
\iqlookfor{renommages et réutilisation~; une correspondance datée vers une clé permanente.}
\end{solution}

\begin{solution}{iq:rs:universe-symbology-and-corporate-actions:2}
\omterm{def:rs:point-in-time-data-and-the-biases:lookahead}{Biais d'anticipation} par les divisions ultérieures~: le cours ajusté à une date ancienne est divisé par les divisions qui ont suivi, survenues après
la hausse des cours, si bien que le plancher écarte les futurs gagnants. Appliquer les règles de niveau au cours échangé ce jour-là.
\iqlookfor{que les facteurs d'ajustement viennent du futur.}
\end{solution}

\begin{solution}{iq:rs:universe-symbology-and-corporate-actions:3}
Une table des titres (\omterm{def:rs:universe-symbology-and-corporate-actions:permanent}{identifiant permanent}, émetteur, dates de cotation et de radiation, motif et \omterm{def:rs:universe-symbology-and-corporate-actions:delisting}{rendement de radiation}) et une table des symboles (type de symbole,
symbole, \omterm{def:rs:universe-symbology-and-corporate-actions:permanent}{identifiant permanent}, date de début, date de fin) indexée sur (symbole, date de début)~; la requête prend la ligne telle que début $\le D <$ fin. Il faut ajouter une
table d'opérations sur titres indexée par \omterm{def:rs:universe-symbology-and-corporate-actions:permanent}{identifiant permanent}, avec les dates d'annonce et les dates ex.
\iqlookfor{lignes à plage de dates et index sur (symbole, début).}
\end{solution}

\begin{solution}{iq:rs:universe-symbology-and-corporate-actions:4}
À partir de la capacité et des coûts~: assez grand pour la largeur (chapitre~15) et assez petit pour que chaque valeur puisse être négociée à la taille de la
stratégie, la liquidité étant mesurée telle qu'elle est connue à chaque date~; un tampon (rang de sortie plus large que la taille) et un historique minimal pour limiter la rotation
et les démarrages à froid~; aucune règle portant sur les niveaux ajustés.
\iqlookfor{liquidité contre largeur, un tampon, des règles à la date.}
\end{solution}

\begin{solution}{iq:rs:universe-symbology-and-corporate-actions:5}
Le rendement correspond à ce qu'un détenteur pourrait réaliser~: le cours de gré à gré après la radiation, ou le recouvrement. S'il manque, une
estimation par motif de radiation (comme le $-30~\%$ que Shumway a retrouvé pour les radiations pour performance), signalée~; ne jamais prendre zéro, ce qui constitue
l'erreur optimiste sur les positions perdantes.
\iqlookfor{le \omterm{def:rs:universe-symbology-and-corporate-actions:delisting}{rendement de radiation}, et ne pas inscrire zéro.}
\end{solution}

\begin{solution}{iq:rs:universe-symbology-and-corporate-actions:6}
Entre deux opérations, tous les cours sont divisés par une même constante, qui s'annule dans un rendement~; à une date ex, le facteur retire la variation mécanique
du cours, si bien que le rendement ajusté égale le rendement de détention. Un plancher sur les cours ajustés à la date $t$ utilise les facteurs de division des
opérations postérieures à $t$, lesquelles sont plus probables pour les actions qui ont ensuite monté.
\iqlookfor{l'argument d'annulation et une fuite concrète.}
\end{solution}
